% generated by R/04_tables.R from table1_mutations_by_fraction.md -- do not edit
% needs \usepackage{booktabs,adjustbox} in the preamble
\begin{table}[htbp]
\centering
\caption{Mutations by sequenced fraction}%
\label{tab:mutfrac}
\begin{adjustbox}{max width=\linewidth}
\begin{tabular}{lrlrrrl}
\toprule
\textbf{Fraction} & \textbf{Clones sequenced} & \textbf{Clones with a mutation} & \textbf{Mutations} & \textbf{Singletons} & \textbf{In $>1$ clone} & \textbf{Highest frequency} \\
\midrule
Early spore fraction & 52 & 2 (4\%) & 2 & 2 & 0 & 1/52 (1.9\%) \\
Endpoint spore fraction & 96 & -- & 4 & 4 & 0 & 1/96 (1.0\%) \\
Endpoint total fraction & 84 & 48 (57\%) & 34 & 24 & 10 & 8/84 (9.5\%) \\
\bottomrule
\end{tabular}
\end{adjustbox}
\par\smallskip
{\footnotesize\raggedright Counts exclude the variant at position 3,529,981, carried by 47 of 84 clones in the endpoint total fraction and therefore not something that arose during the experiment. No clone-by-variant matrix survives for the endpoint spore fraction, so its per-clone column is blank; its four mutations are known to be singletons.\par}
\end{table}
